Effective knowledge transfer across bioprocessing systems is essential for reducing process development timelines, adapting platforms to new cell lines and products, facilitating scale-up and optimising production \cite{Reichelt2017BioprocessCorrelations, Helleckes2023MachinePractice}. However, inherent biological variability and system-specific differences in metabolism, growth dynamics, and, consequently, nutrient uptake rates often limit the transferability of process insights across systems \cite{Chen2019AnCells}. In the absence of a structured methodology for adopting prior knowledge, process development continues to rely on extensive experimentation, leading to increased cost and longer timelines \cite{Kroll2017Model-BasedLifecycle}.

Mathematical models offer a systematic framework for understanding and predicting bioprocess behaviour through quantitative description of system dynamics \cite{Kiparissides2011ClosingVitro}. Mechanistic models, in particular, are widely employed due to their ability to explain biological and physicochemical behaviours based on first principles \cite{Kyriakopoulos2018KineticBiomanufacturing, Almquist2014KineticPerformance}. While these models can provide valuable insight into cell behaviour, their predictive performance is often constrained by system-specific parameterisation \cite{Bayer2023ComparisonPaths}. When applied to different bioreactor scales, cell lines, or operating regimes, mechanistic models typically require extensive reparametrisation to maintain accuracy \cite{Sha2018MechanisticProduction}.
 
To overcome these challenges, recent efforts have focused on data-driven and hybrid modelling approaches designed to improve model transferability. Stand-alone data-driven methods, such as transfer learning using Artificial Neural Networks (ANNs) \cite{Rogers2022AData} or multi-fidelity Gaussian process models \cite{Sun2022Multi-fidelityData, Martens2025HolisticOptimization}, are flexible and often perform well in data-scarce regimes \cite{KarimiAlavijeh2024ADevelopments}. However, they typically lack mechanistic interpretability. Hybrid models improve generalizability while partially retaining interpretability by combining mechanistic structure with data-driven components, as demonstrated in several recent studies \cite{Narayanan2019AProteins, Bayer2021ModelExperiments, Bayer2023ComparisonPaths, Pinto2023HybridNetworks, Ramos2024DeepEquations, Cruz-Bournazou2022HybridCultures}. Yet in many cases, the mechanistic backbone is reduced to mass balances, with core biological behaviours such as metabolic fluxes and growth kinetics captured by black-box surrogates \cite{Narayanan2019AProteins, Tsopanoglou2021MovingBioprocesses}, reducing the model's capacity to provide mechanistic insight. In other hybrid implementations, system identity is encoded via one-hot vectors \cite{Helleckes2024NovelBioprocesses} or embedding layers \cite{Hutter2021KnowledgeVectors}, allowing the model to distinguish between systems based on historical data. However, they do not directly identify the underlying physiological traits responsible for these differences. Moreover, although hybrid models are less data-intensive than fully data-driven approaches, they still require multiple datasets to train their data-driven components \cite{Bayer2023ComparisonPaths}. This presents a challenge in early-stage process development, where reliable predictive models are most needed but data is scarce \cite{Lu2024Multi-cell-lineModels}. In addition, the absence of explicit uncertainty handling increases the risk of overfitting, particularly when datasets are limited or noisy \cite{Schweidtmann2024AMethodologies}. 

A fundamental limitation shared across mechanistic, data-driven, and hybrid models is their reliance on fixed parameter values across process duration. Once parametrised, these models operate deterministically and are typically applied offline, limiting their ability to adapt to deviations in system behaviour caused by parameter uncertainty, measurement noise, or biological variability \cite{HernandezRodriguez2021DynamicTrain, Sinner2022OnlineMixture}. This motivates the need for dynamic modelling frameworks that incorporate real-time measurements to update both states and parameters while retaining mechanistic structure. In this context, state estimation algorithms offer a practical alternative.

Various state estimation techniques have been applied to bioprocess applications, each with distinct advantages and limitations. The Extended Kalman Filter (EKF) remains one of the most widely used methods due to its easy implementation and computational efficiency \cite{Ohadi2015DevelopmentCultures, Kramer2016On-lineCerevisiae, Gopalakrishnan2011IncorporatingEstimation, Narayanan2020Hybrid-EKF:Culture, 
Dewasme2019StateValidation, Bogaerts2004ParameterSensors, Ulonska2018ModelFed-batch}. EKF linearises the process model and is generally suitable for mildly nonlinear systems \cite{Lyubenova2021Model-basedReview}. However, its accuracy deteriorates when applied to strongly nonlinear and discontinuous processes, which are common in mammalian cell cultures \cite{Yousefi-Darani2021TheProcesses}. The Unscented Kalman Filter (UKF) addresses nonlinearities more effectively than the EKF by propagating sigma points through the nonlinear model \cite{Wan2000UnscentedEstimation}, with successful applications reported in the bioprocessing literature \cite{Fernandes2015ExtendedCultures, Dewasme2015StateCells, Abadli2021AnCultures, Kramer2019AFilter, Kolas2009ConstrainedApproach, Yousefi-Darani2021ParameterFilter, Paquet-Durand2022Online, Kemmer2023NonlinearCultivations}. While more robust than the EKF \cite{Wan2000UnscentedEstimation}, the performance of UKF is sensitive to the choice of hyperparameters, often requiring careful tuning to ensure stability and accuracy \cite{Dunik2012UnscentedScaling, Scardua2017CompleteTuning}. Particle Filter (PF) has been applied to many complex bioprocess systems that exhibit highly nonlinear dynamics and non-Gaussian uncertainty \cite{Golabgir2016CombiningProduction, Stelzer2017ComparisonBioprocesses, Kager2019ExtensionMeasurements, Sinner2022OnlineMixture, Kager2018StateMeasurements}, at the expense of higher computational demand \cite{Elfring2021ParticleTutorial}. Apart from recursive methods, Moving Horizon Estimation (MHE) formulates the state estimation problem as a constrained optimisation over a finite time horizon \cite{Robertson1996AEstimation, Rao2001ConstrainedApproach, Haseltine2005CriticalEstimation}, explicitly accounts for model constraints and past measurements \cite{Tuveri2022BioprocessApplication, HernandezRodriguez2021DynamicTrain, Alexander2020ChallengesProcesses}. However, the need to solve an optimisation problem introduces substantial computational overhead, which can limit its use for real-time applications.

The Ensemble Kalman Filter (EnKF) offers a promising solution to many of the challenges associated with state estimation techniques of bioprocessing systems. Originally developed for large-scale systems in geosciences and reservoir modelling \cite{Evensen2022DataProblem}, EnKF combines the computational efficiency of Kalman-based methods with the flexibility of ensemble sampling to handle nonlinearity and uncertainty. Unlike the EKF and UKF, it requires no Jacobian computation or sigma point selection, avoids the high computational burden of PF,  and the need to solve an optimisation problem in MHE. By propagating an ensemble of model realisations, the EnKF naturally captures the evolution of uncertainty in both states and parameters, without requiring the explicit computation or storage of process covariance matrices \cite{Evensen2009TheSystems}. This ensemble-based formulation offers a favourable balance among estimation accuracy, ease of implementation, and computational efficiency. 

We have previously applied the EnKF for soft sensing of unmeasured intracellular states in CHO cell cultures \cite{Yu2023EnsembleAntibodies}. In this study, we extend its application beyond state estimation to the recursive updating of uncertain kinetic parameters throughout the duration of the culture. By analysing the temporal evolution of parameter ensembles, we could gain insights into dynamic changes in cellular behaviour that would otherwise be masked by static calibration. As the parameter trajectories are continuously refined based on incoming measurements, this approach preserves the mechanistic structure while adapting flexibly to new conditions without the need for reparametrisation. This capability supports knowledge transfer by allowing a model developed for one specific system to be applied to different cell lines, scales, or operating conditions using only a single experimental dataset. Moreover, the evolving parameter trajectories enable a form of dynamic sensitivity analysis, where the propagation of ensemble uncertainty highlights which parameters are most influential under the new system conditions.

